\chapter{Introduction}\label{chap:introduction}

Tout matériau réel est imparfait : un atome manque, un autre occupe la mauvaise place ou une impureté s'est immiscée. À l'échelle macroscopique, ces accidents se comptent par milliards et déterminent souvent ce que le matériau fait de mieux comme de pire : ils limitent la mobilité des électrons dans un transistor, déterminent la couleur de certains cristaux, ancrent les vortex d'un supraconducteur ou font apparaître un moment magnétique là où le matériau parfait n'en possède pas. Comprendre ces défauts, leur nature, leur concentration et surtout leur influence sur les propriétés du matériau, constitue un enjeu central de la science des matériaux.

\section{Contexte et motivation}
\label{sec:contexte}

L'étude des défauts prend une importance particulière dans les matériaux bidimensionnels. Lorsqu'un matériau est réduit à une seule couche d'atomes, chaque atome est un atome de surface et chaque défaut est directement accessible à l'expérience. Ces matériaux, dont le plus connu est le graphène, sont ainsi un terrain idéal pour étudier les défauts et leurs effets sur les propriétés macroscopiques.

\bloc{Le graphène}

Le graphène est une feuille de carbone d'un seul atome d'épaisseur, dont les atomes sont arrangés sur un réseau en nid d'abeille. Isolé pour la première fois en 2004 par l'exfoliation mécanique du
graphite~\cite{novoselovElectricFieldEffect2004}, il est à la fois l'un des matériaux les plus simples que l'on puisse imaginer et l'un des plus singuliers~\cite{geimRiseGraphene2007}. Sa cellule unitaire ne contient que deux atomes. Pourtant, avec un module de Young de l'ordre du térapascal et une résistance à la rupture de $\sim 130$~GPa, la plus haute jamais mesurée~\cite{leeMeasurementElasticProperties2008}, c'est le matériau le plus résistant que l'on connaisse. Il conduit la chaleur mieux que le cuivre, à $\sim 5 \times 10^{3}$~W/mK~\cite{balandinSuperiorThermalConductivity2008}, et les électrons avec une mobilité qui dépasse $10^{5}$~cm$^{2}$\,V$^{-1}$\,s$^{-1}$~\cite{bolotinUltrahighElectronMobility2008}. En outre, ce matériau, qui bat tous les records, est transparent à $97.7$~\% dans le visible~\cite{nairFineStructureConstant2008}.

Sa richesse ne tient pas à la complexité de sa chimie, mais à sa géométrie et à la nature des liens entre ses atomes. Le réseau en nid d'abeille mène à une structure électronique particulière : la bande de valence $\pi$ touche linéairement la bande de conduction $\pi^*$ aux points de Dirac $K$ et $K'$. À ces points, le gap et la densité d'états s'annulent : le graphène est à l'intersection d'un métal et d'un semi-conducteur, d'où sa catégorisation en tant que \textit{semi-métal}. Ses électrons de basse énergie se comportent comme des particules relativistes sans masse~\cite{castronetoElectronicPropertiesGraphene2009}. Cette structure électronique est à l'origine d'une grande partie de ses propriétés emblématiques et, nous le verrons, de la richesse de sa réponse à la lumière.

Ces propriétés font du graphène un système modèle pour l'étude des défauts pour deux raisons. D'abord, le graphène parfait est très étudié et exceptionnellement bien compris, de sorte que l'effet d'un défaut se mesure par rapport à une référence connue et solide et peut être quantifié efficacement. Ensuite, ses singularités en font un cas exigeant : les approximations physiques et les méthodes de calcul qui conviennent à un métal ou un semi-conducteur sont souvent mises à l'épreuve. Ainsi, une méthode qui traite correctement le graphène a de bonnes chances de se généraliser à des matériaux moins exigeants.

\bloc{Les défauts}

Les défauts ne sont pas qu'une curiosité théorique : ils sont présents dans tous les échantillons, qu'on le veuille ou non. La croissance chimique en phase vapeur, le transfert d'un substrat à l'autre, l'exposition à un faisceau d'ions ou d'électrons en introduisent tous dans le graphène~\cite{banhartStructuralDefectsGraphene2011}. Les défauts sont aussi souvent voulus, car ils peuvent conférer au matériau des propriétés nouvelles : greffer de l'hydrogène au graphène lui ouvre un gap~\cite{eliasControlGraphenesProperties2009} ; des lacunes ou des atomes étrangers fixent plus fortement les molécules adsorbées~\cite{zhangImprovingGasSensing2009}, ce qui améliore des capteurs déjà capables de détecter une molécule unique~\cite{schedinDetectionIndividualGas2007} ; et des pores de quelques nanomètres font du graphène une membrane à travers laquelle des molécules d'ADN peuvent être détectées une à une~\cite{garajGrapheneSubnanometreTranselectrode2010}.

Qu'ils soient accidentels ou introduits délibérément, les défauts peuvent modifier profondément les propriétés physiques et chimiques du matériau. Pour quantifier ce changement et l'attribuer à un type de défaut donné, il faut une sonde.

\bloc{La spectroscopie Raman}

La lumière est une sonde de choix : éclairer un cristal avec un laser et regarder la lumière qui en ressort est, en apparence, la plus simple des expériences. Presque toute la lumière diffusée a la couleur du laser, sauf une infime fraction qui en ressort décalée en fréquence : elle a laissé dans le matériau un quantum de vibration, un phonon, et ce décalage en mesure l'énergie. C'est l'effet Raman~\cite{ramanNewTypeSecondary1928}, et le spectre de ces décalages est une véritable empreinte digitale du réseau cristallin du matériau.

Pour le graphène, cette empreinte est devenue un outil de caractérisation incontournable~\cite{ferrariRamanSpectrumGraphene2006,ferrariRamanSpectroscopyVersatile2013a}. La mesure est rapide, non destructive, se fait à l'air ambiant sans aucune préparation complexe de l'échantillon et permet, sur un seul spectre, d'en apprendre sur le nombre de couches, la déformation (\textit{strain})~\cite{mohiuddinUniaxialStrainGraphene2009}, le dopage et le désordre. Le spectre du graphène parfait est dominé par deux pics : le pic G et le pic 2D. Le premier est la signature ordinaire d'un phonon dans le matériau, tandis que le second ne l'est pas : il fait intervenir deux phonons à la fois. Dans la plupart des matériaux, un tel processus est beaucoup moins intense que celui à un seul phonon. Dans le graphène, il est pourtant souvent plus intense. La raison tient à la structure électronique décrite précédemment : en l'absence de gap, tout laser visible peut induire des transitions électroniques réelles, ce qui amplifie fortement la diffusion, et davantage encore pour le pic 2D, qui est causé par deux transitions électroniques réelles. Nous parlons alors de diffusion Raman \textit{résonante} pour le pic G et \textit{doublement résonante} pour le pic 2D.

Cette résonance est à la fois ce qui contribue à la richesse du spectre Raman du graphène et ce qui complique son calcul théorique. Pour décrire correctement le spectre résonant, il faut considérer le processus à l'échelle de ses constituants : les quanta de lumière, les électrons, les phonons, les défauts, ainsi que les interactions entre eux. À cette échelle, une description quantique explicite du processus est nécessaire.

Les défauts laissent eux aussi leur signature dans ce spectre. Tuinstra et Koenig ont observé dans le spectre Raman du graphite un pic dont l'intensité croît avec le désordre~\cite{tuinstraRamanSpectrumGraphite1970} et est absent du graphite parfait : le pic D. \textit{A priori}, ce pic est une énigme car un défaut statique ne peut ni absorber ni céder d'énergie, et ne peut donc pas, à lui seul, décaler la fréquence de la lumière. L'explication est microscopique : le défaut n'agit pas sur la lumière, mais sur les électrons qu'elle a excités et leur permet d'émettre un phonon qui leur serait autrement interdit.

Aujourd'hui, l'intensité du pic D relative à celle du pic G est devenue la mesure de routine de la densité de défauts dans le graphène~\cite{luccheseQuantifyingIoninducedDefects2010,martinsferreiraEvolutionRamanSpectra2010}. Néanmoins, un spectre n'est utile que s'il est comparé à une référence et l'expérience ne peut fournir cette référence que pour les rares défauts que nous pouvons produire isolément et en concentration contrôlée. Pour tout le reste, le rôle revient à la théorie et à la simulation numérique : prédire, pour chaque type de défaut et chaque concentration, le spectre Raman attendu, à partir de rien d'autre que la structure du matériau et la nature du défaut. C'est ce que promettent les calculs à partir des premiers principes, dits \textit{ab initio} : calculer les propriétés d'un matériau, y compris sa réponse Raman, à partir des lois de la mécanique quantique, sans ajuster les paramètres d'un modèle sur l'expérience. Ce mémoire pose, en partie, les fondations d'un tel calcul.

\newpage
\bloc{Les calculs à partir des premiers principes}

La théorie du spectre Raman résonant du graphène avec des défauts existe déjà~\cite{venezuelaTheoryDoubleresonantRaman2011}, mais les descriptions existantes reposent généralement sur des modèles paramétrés. De telles approches permettent de reproduire les caractéristiques essentielles du spectre, mais ne peuvent pas servir de référence indépendante aux expérimentateurs puisque les paramètres du modèle eux-mêmes doivent être ajustés, pour chaque défaut et chaque matériau, sur l'expérience ou des calculs plus complets.

Nous voulons prédire plutôt que de reproduire : décrire le matériau à partir de sa structure et de ses constituants microscopiques, à l'aide des seules lois de la mécanique quantique. C'est ce que permettent la théorie de la fonctionnelle de la densité (DFT)~\cite{hohenbergInhomogeneousElectronGas1964,kohnSelfConsistentEquationsIncluding1965} et sa version perturbative. Introduire un défaut revient alors à modifier cette structure et refaire le calcul : son effet sur les propriétés en découle directement sans paramètre choisi à l'avance. Le même calcul s'applique ensuite, en principe, à n'importe quel type de défaut et pour n'importe quel matériau. Il suffit de fournir une autre structure d'entrée, toujours sans ajuster aucun paramètre.

Cette approche a un second avantage. Un spectre Raman résonant se calcule à partir d'un ensemble d'ingrédients et dans un modèle, chacun est ajusté séparément sans garantit qu'ils soient compatibles entre eux. À l'inverse, s'ils sont issus d'un même calcul DFT, ils le sont par construction.

À la lumière\footnote{C'est le cas de le dire !} de ces avantages, une question s'impose : pouvons-nous obtenir ces ingrédients à partir des premiers principes pour le graphène avec des défauts, de sorte que le spectre Raman puisse ensuite être calculé sans aucun paramètre physique ajusté ? C'est la question qui guide ce mémoire.

\section{État de l'art et positionnement}
\label{sec:intro-etat-art}

Deux autres questions situent ce mémoire dans la littérature : jusqu'où le spectre Raman résonant des matériaux avec défauts se calcule-t-il aujourd'hui à partir des premiers principes, et quels sont les obstacles principaux ?

\bloc{Le spectre Raman résonant à partir des premiers principes}

Pour un cristal parfait, loin de toute résonance, le spectre Raman se calcule de façon routinière~\cite{lazzeriFirstPrinciplesCalculationVibrational2003} et les principaux codes DFT l'implémentent. Toutefois, l'approche employée par ces codes suppose que le laser incident dans l'expérience Raman ne crée aucune excitation réelle dans le matériau. Cette hypothèse est systématiquement violée à toutes les énergies de laser dans le graphène en raison de son gap nul, et cette approche n'est pas applicable à notre cas.

À l'inverse, les calculs de spectres résonants sont nettement plus rares. Ils requièrent une théorie quantique et imposent des exigences numériques draconiennes, car l'intensité est dominée par les quelques états électroniques qui satisfont exactement la condition de résonance. Gillet, Giantomassi et Gonze~\cite{gilletFirstprinciplesStudyExcitonic2013a} ont établi un tel cadre au premier ordre en incluant les effets excitoniques. Reichardt et Wirtz ont calculé l'intensité du pic G du graphène \textit{sans} défaut à partir d'ingrédients tous issus des premiers principes~\cite{reichardtInitioCalculation$G$2017}, puis ont formulé une théorie plus générale~\cite{reichardtTheoryResonantRaman2019}. Le spectre Raman résonant du graphène \textit{parfait} est donc accessible aujourd'hui \textit{ab initio}.

Néanmoins, celui du graphène avec des défauts ne l'est toujours pas. La théorie complète du Raman résonant en présence de défauts a été formulée par Venezuela, Lazzeri et Mauri~\cite{venezuelaTheoryDoubleresonantRaman2011}, qui ont reproduit les pics D et 2D ainsi que leur dispersion avec l'énergie du laser. Leurs défauts sont toutefois modélisés en étant introduits dans un modèle de liaisons fortes où leur couplage aux électrons dépend de paramètres choisis plutôt que calculés. Même le code d'intensité Raman résonante le plus récent suppose un couplage électron-défaut constant~\cite{huangQR$^2$codeOpensourceProgram2026}. À notre connaissance, aucun spectre Raman résonant avec des défauts n'a été calculé, pour quelque matériau que ce soit, avec tous ses ingrédients issus des premiers principes. Le maillon manquant est le couplage \textit{ab initio} électron-défaut et ce mémoire le construit.

\bloc{Le couplage électron-défaut à partir des premiers principes}

Le calcul \textit{ab initio} du couplage entre les électrons et un défaut n'est pas une opération de coutume en science des matériaux. Aucun code communautaire n'en fait une quantité de routine et la littérature regroupe quelques articles basés sur des implémentations spécialisées. Le contraste avec le couplage électron-phonon est frappant : ce dernier dispose d'un cadre théorique établi et validé~\cite{giustinoElectronphononInteractionsFirst2017,giustinoElectronphononInteractionUsing2007} et les codes qui l'implémentent sont distribués à la communauté depuis près de deux décennies~\cite{ponceEPWElectronphononCoupling2016,leeElectronPhononPhysics2023a,zhouPerturboSoftwarePackage2021,gonzeAbinitProjectImpact2020,bruninPhononlimitedElectronMobility2020}. La raison principale est que la présence d'un défaut brise explicitement la symétrie de translation du système, et que le moyen usuel est de contourner ce problème est de recourir à des super-cellules coûteuses.

Deux approches récentes ont posé les bases et chacune apporte la moitié de la solution. Lu, Zhou et Bernardi~\cite{luEfficientInitioCalculations2019,luInitioElectrondefectInteractions2020} ont montré comment calculer le couplage dans une base d'ondes planes à partir du potentiel d'une super-cellule contenant le défaut, puis comment l'interpoler par fonctions de Wannier sur les grilles denses qu'exigent les propriétés de transport ou spectrales, comme l'intensité Raman. Leur formalisme traite toutefois la diffusion au premier ordre dans le potentiel du défaut, l'approximation de Born, ce qui convient aux impuretés chargées des semi-conducteurs, mais échoue pour les défauts dont le potentiel est intense, comme la lacune. Kaasbjerg~\cite{kaasbjergAtomisticMatrixTheory2020} a au contraire traité la diffusion à tous les ordres et a montré que l'approximation de Born peut se tromper d'un ordre de grandeur pour les défauts de matériaux bidimensionnels. Son calcul est toutefois fait dans une base d'orbitales atomiques localisées, sans interpolation vers les grilles denses.

Le spectre Raman résonant a besoin des deux à la fois : un traitement exact du potentiel de défaut \textit{et} une évaluation sur des grilles de plusieurs centaines de milliers de points pour satisfaire ses exigences numériques intransigeantes. Combiner les deux n'est pas immédiat et soulève plusieurs problèmes. D'une part, l'interpolation par fonctions de Wannier exige que la grille de la super-cellule et celle des fonctions de Wannier soient commensurables, ce qui forcerait des super-cellules de plus d'un millier d'atomes et est complètement intraitable. D'autre part, traiter la diffusion à tous les ordres sur les grilles requises par le Raman requiert de manipuler des matrices dont la dimension est de plusieurs centaines de milliers, ce qu'aucune machine ne peut faire. Ce mémoire résout ces deux problèmes et livre, pour une lacune dans le graphène, le couplage électron-défaut \textit{ab initio} interpolé par fonctions de Wannier sur des grilles arbitrairement denses et la diffusion qu'il induit, traitée à tous les ordres dans le potentiel de défaut.

\section{Contributions et plan du mémoire}
\label{sec:intro-contributions}

Les deux chapitres suivants posent le cadre du mémoire. Le chapitre~\ref{chap:théorie} détaille les outils : la théorie de la fonctionnelle de la densité pour les électrons, sa version perturbative pour les phonons, les fonctions de Wannier pour l'interpolation vers les grilles denses qu'exige la résonance et l'Hamiltonien complet du système sur lequel tout le reste est construit. Le chapitre~\ref{chap:Raman} rassemble ces outils pour établir la théorie de la diffusion Raman, de sa description semi-classique aux amplitudes quantiques du premier et du deuxième ordre, y compris les processus faisant intervenir des défauts. Il se clôt sur la liste des ingrédients que ces amplitudes requièrent : les énergies des électrons $\varepsilon_{n\kv}$ et des phonons $\omega_{\nu\qv}$, les couplages électron-photon $g^{(\lambda)}$, électron-phonon $g^{(\nu)}$ et électron-défaut $M$, ainsi que les taux d'amortissement $\Gamma^\text{ep}_{n\kv}$ et $\Gamma^\text{ed}_{n\kv}$ et la largeur des phonons $\gamma_{\nu\qv}$. Le chapitre~\ref{chap:théorie} fournit les énergies et définit le couplage électron-photon, dont l'évaluation numérique est laissée hors du cadre de ce mémoire. Les deux chapitres suivants calculent tous les autres ingrédients.

Le chapitre~\ref{chap:défauts} est le c\oe ur du mémoire et contient sa contribution originale : le couplage électron-défaut est construit à partir des premiers principes. Nous calculons d'abord les éléments de matrice $M$ dans une base d'ondes planes en utilisant un raffinement du potentiel de défaut par remplissage par zéros, exact pour un potentiel localisé, qui découple la taille de la super-cellule de la grille des fonctions de Wannier, ouvrant la voie à l'interpolation par ces dernières. Nous employons ensuite un schéma d'amas, adapté de Kaasbjerg~\cite{kaasbjergAtomisticMatrixTheory2020} à la base de Wannier, qui permet de traiter la diffusion électron-défaut à tous les ordres sans manipuler des objets de très grande dimension pour obtenir le taux d'amortissement électronique dû aux défauts $\Gamma_{n\kv}^\text{ed}$. Nous appliquons notre méthode à une lacune dans le graphène, non relaxée, sans polarisation de spin et pour une faible concentration de défauts, l'objectif étant de valider la méthode plutôt que de décrire la lacune réelle. Elle livre trois résultats : un état quasi-lié sous la singularité de van Hove de la bande $\pi$, l'échec de l'approximation de Born qui surestime le taux d'un facteur trois en médiane et diverge aux singularités de van Hove, et un couplage dominé par sa partie non locale.

Le chapitre~\ref{chap:eph} calcule le couplage électron-phonon $g^{(\nu)}$, le taux $\Gamma^\text{ep}_{n\kv}$ et la largeur $\gamma_{\nu\qv}$, avec le code \texttt{EPW}, appliqué au graphène parfait. Il est volontairement plus court : la méthode est déjà établie et le chapitre se concentre sur la validation des résultats. Il se termine par la comparaison des deux mécanismes d'amortissement électronique : phonons et défauts. Nous trouvons qu'ils sont du même ordre de grandeur à température ambiante et pour une concentration de défauts de l'ordre du pour cent, de sorte qu'aucun des deux mécanismes n'est négligeable.

Pour finir, le titre du mémoire commence par \og Vers \fg{} et ce n'est pas une faute de frappe : le spectre Raman résonant du graphène avec des défauts, dans toute sa gloire, n'est pas calculé ici. Seuls ses ingrédients, à un près, le sont ; ils sont validés, puis livrés sur les grilles denses que le spectre résonant exige. Outre le couplage électron-photon, trois étapes séparent encore ces ingrédients d'un calcul d'intensité en bonne et due forme. D'abord, les valeurs propres de Kohn-Sham ne sont pas des énergies d'excitation et la fonctionnelle que nous utilisons sous-estime le couplage électron-phonon  et électron-défaut dans le graphène~\cite{lazzeriImpactElectronelectronCorrelation2008} : tous requièrent une correction par l'approximation $GW$, qui traite la corrélation électronique au-delà de la DFT en combinant la fonction de Green des électrons $G$ et l'interaction de Coulomb écrantée $W$~\cite{hedinNewMethodCalculating1965}. Ensuite, la super-cellule doit être relaxée pour capturer la reconstruction de la lacune dans le graphène~\cite{el-barbaryStructureEnergeticsVacancy2003,banhartStructuralDefectsGraphene2011}. Enfin, le calcul de super-cellule doit être polarisé en spin, puisque la lacune induit un moment magnétique~\cite{yazyevDefectinducedMagnetismGraphene2007}. Les deux premièrent étapes relèvent du calcul et la troisième demande d'étendre le formalisme au spin. Nous les laissons à des études doctorales. La méthode, soit le calcul du couplage électron-défaut \textit{ab initio} et son insertion dans la théorie du Raman résonant, est ce que ce mémoire établit.